% =============================================================================
% RESUMEN (hasta 250 palabras) / ABSTRACT (up to 250 words)
% =============================================================================
\chapter*{Resumen}
\addcontentsline{toc}{chapter}{Resumen}

% 
TODO: completar con el resumen real de la tesis
%Los procesos de acreditación de carreras universitarias ante la Comisión Nacional de Evaluación y Acreditación Universitaria (CONEAU) demandan una elevada carga administrativa y técnica, que implica gestionar cientos de evidencias institucionales, coordinar múltiples actores y garantizar la trazabilidad de la documentación presentada. Esta complejidad genera ineficiencias, duplicaciones de esfuerzo y baja reproducibilidad en los informes institucionales.

%Esta tesis propone, diseña e implementa un \textbf{Sistema Multiagente Inteligente} que integra \textbf{Modelos de Lenguaje de Gran Escala (LLM)} y el \textbf{Model Context Protocol (MCP)} como capa de interoperabilidad, con el fin de automatizar y optimizar el seguimiento de los procesos de acreditación ante la CONEAU. El sistema permite la extracción automática de evidencias desde documentos institucionales, su clasificación semántica según los estándares normativos vigentes, y la generación de reportes explicativos que asisten a las comisiones de autoevaluación.

%La arquitectura implementada comprende un backend REST desarrollado con FastAPI (Python), un frontend web con React, y un módulo de indexación vectorial para la búsqueda semántica de propuestas docentes. Los resultados preliminares demuestran que el sistema reduce significativamente el tiempo requerido para la recolección y clasificación de evidencias, mejorando la consistencia y trazabilidad del proceso.

\bigskip
\noindent\textbf{Palabras clave:} completar...%acreditación universitaria, CONEAU, sistemas multiagente, modelos de lenguaje, MCP, extracción de información, búsqueda semántica, FastAPI, React.

\bigskip
\hrule
\bigskip

\chapter*{Abstract}
\addcontentsline{toc}{chapter}{Abstract}

% 
TODO: completar con el resumen real de la tesis
%University accreditation processes before Argentina's National Commission for University Evaluation and Accreditation (CONEAU) require extensive administrative and technical effort, involving the management of hundreds of institutional evidences, coordination of multiple stakeholders, and ensuring documentary traceability. This complexity leads to inefficiencies, duplicated effort, and low reproducibility in institutional reports.

%This thesis proposes, designs, and implements an \textbf{Intelligent Multi-Agent System} that integrates \textbf{Large Language Models (LLMs)} and the \textbf{Model Context Protocol (MCP)} as an interoperability layer, aiming to automate and optimize the monitoring of CONEAU accreditation processes. The system enables automatic extraction of evidences from institutional documents, their semantic classification against current regulatory standards, and the generation of explanatory reports that assist self-assessment committees.

\bigskip
\noindent\textbf{Keywords:} completar %university accreditation, CONEAU, multi-agent systems, large language models, MCP, information extraction, semantic search, FastAPI, React.
